本文针对题目提出的模型综合评价、场景选择与成本约束决策问题，构建了由基础能力、场景适配和预算约束组成的递进模型。

问题一在 \(164/210=78.1\%\) 的有效公开观测下，先区分普通缺失与结构性缺失，再通过缺失感知相关性、Family 平衡成对比较和正则化 Bradley--Terry 得到五维能力结构。Ranking A 的点估计中，GPT-5.6 Sol、Claude Fable 5 和 Kimi K3 位列前三；但 Bootstrap 区间较宽，Kimi K3 的 Top-3 经验概率为 \(0.4225\)，因此该排序应理解为正式数据下的中心评价，而不是绝对稳定的市场名次。

问题二说明综合能力并不等于任务适配。科研场景中 Claude Fable 5 居首，原因是 C1--C3 的能力组合较均衡，在负替代弹性下短板惩罚较小；代码场景中 Kimi K3 居首，原因是代码工程能力 C4 在场景权重中占主导；GPT-5.6 Sol 虽然综合排名第一，但在科研和代码场景中的位置发生变化。这说明模型选择必须经过能力需求映射，而不能停留在总体榜单。

问题三进一步将场景效用与工作负载成本结合。代码场景的预算决策依次经过 DeepSeek-V4-Flash、DeepSeek-V4-Pro 和 Kimi K3，切换阈值为 \(0.01056\)、\(0.03168\) 和 \(0.10683\) USD。效用优势只有在成本信息完整且预算允许时，才能转化为可执行的部署方案；因此，效用高但成本不可核验的模型不应被强行纳入严格 Pareto。

综上，题目的最终决策逻辑为
\[
\boxed{
\text{基础能力识别}
\longrightarrow
\text{场景适配}
\longrightarrow
\text{预算约束下的决策}
}.
\]
若关注综合能力，应参考 Ranking A 的中心评价并同时查看不确定性；若属于科研任务，应优先比较科研场景中的效用与短板结构；若属于代码任务，应重点考察 C4 主导下的场景效用；若成本敏感，则沿相应场景的 Pareto 前沿选择；随着预算提高，再根据预算决策区间带切换到更高效用节点。由此，不存在脱离任务需求与预算约束的绝对最优模型。
